 
\section{Problem Formulation}
\label{sec:problem}

\subsection{Inputs and Objective}

We assume a fixed set of high-level safety requirements $P$, supplied by
the developer before agent evolution. For example, $P$ may require that
confidential information be accessible only to authorized parties.
Our method does not infer these requirements from agent updates. Instead,
it identifies additional executable obligations needed to preserve them
as the agent's capabilities evolve.

Let
\[
A_{t-1}=\langle M_{t-1},R_{t-1},T_{t-1},W_{t-1}\rangle
\]
denote the deployed agent, comprising its model, memory, tools, and
workflow. An evolution step produces a candidate $\widetilde A_t$
and a structured delta
\[
\Delta A_t=\operatorname{Diff}(A_{t-1},\widetilde A_t).
\]
The inputs to contract evolution are $P$, the current contract
$C_{t-1}$, the candidate agent and its delta, and available environment
evidence. The policy determines which consequences are safety-relevant,
while the delta directs the analysis toward changes that may produce
those consequences.

\paragraph{Contract representation.}
We represent contracts using the Expert Specification Language
(ESL) introduced by RvLLM~\cite{zhang2025rvllm}.
An ESL specification consists of variables, domain-specific predicates,
and logical rules. Each predicate is associated with a natural-language
description defining its intended semantics. For example,
$\operatorname{Authorized}(r,f)$ may denote that recipient $r$ is
authorized to access file $f$.
Rules express relationships among predicates using a restricted
implication form, with universal quantification implicit. For example,
\[
\operatorname{Share}(f,r)
\Rightarrow
\operatorname{Authorized}(r,f)
\]
requires that a file-sharing operation target an authorized recipient.

ESL rules are interpreted against concrete execution contexts.
Following RvLLM, relevant objects and predicate assignments are
extracted from the available context, and rule variables are bound
to these objects to obtain grounded propositions. Predicate truth
values are determined from contextual evidence using three-valued
semantics: $\{\mathsf{True},\mathsf{False},\mathsf{Unknown}\}$.
The last value represents insufficient evidence under an open-world
assumption, rather than a definite violation or satisfaction.
This separates the declarative specification of an obligation
from its context-dependent interpretation.
 
\paragraph{Example: grounding ESL rules in agent execution.}
Consider an agent assisting with document distribution. Its existing
contract requires files to be sent only to authorized recipients:
\[
c_1:\quad
\operatorname{Send}(f,r)
\Rightarrow \operatorname{Authorized}(r,f).
\]
The predicates are defined in natural language:
$\operatorname{Send}(f,r)$ means that the agent sends file $f$
directly to recipient $r$, while $\operatorname{Authorized}(r,f)$
means that $r$ is permitted to access $f$.

Suppose the agent invokes
\texttt{send\_file("report.pdf", "board@company.com")}.
From the tool call and environment context, the interpreter
identifies the file and recipient, binding
$f=\texttt{report.pdf}$ and
$r=\texttt{board@company.com}$.
It then evaluates the grounded predicates. If the recipient
is authorized, both predicates evaluate to $\mathsf{True}$,
and the rule is satisfied. If authorization cannot be
established, the corresponding predicate evaluates to
$\mathsf{Unknown}$ rather than being assumed true.

Now suppose the agent evolves by introducing a reusable
\texttt{upload\_and\_share} tool that uploads files to a
publicly accessible location and returns a link.
When the agent uses this tool for a confidential report,
the original rule $c_1$ does not constrain the upload:
no direct $\operatorname{Send}(f,r)$ operation occurs.
Consequently, the existing contract can remain satisfied
even though the public link violates the fixed policy
that confidential files must be accessible only to
authorized parties.

This execution provides a missing-obligation witness.
Contract evolution can address the gap by introducing
predicates for public accessibility and confidentiality,
together with a new rule:
\[
c_2:\quad
\operatorname{Upload}(f,l)
\land \operatorname{Confidential}(f)
\Rightarrow \neg\operatorname{Public}(l).
\]
Here, $\operatorname{Public}(l)$ is defined as ``location
$l$ is accessible without recipient-specific authorization.''
The new rule constrains the previously uncovered consequence
without prohibiting legitimate uploads of public files. 


\paragraph{Contract evolution.}
A contract $C_t$ contains the ESL predicate definitions and logical
rules applicable to the evolved agent. Clause identifiers and
justifications are maintained as metadata. We distinguish contract
evolution from changes to observation adapters and runtime
instrumentation, denoted by $B_t$:
\[
\underbrace{C_{t-1}\rightarrow C_t}_{\text{new or refined obligations}}
\qquad
\underbrace{B_{t-1}\rightarrow B_t}_{\text{observation and enforcement}}.
\]
For example, adding instrumentation to observe a previously
unavailable predicate changes $B_t$, whereas introducing a rule
that constrains a newly exposed sharing operation changes $C_t$.
Updating $B_t$ alone is therefore not contract co-evolution.

Our objective is to evolve $C_{t-1}$ into $C_t$ such that the
deployed agent $A_t$ satisfies the updated obligations while
preserving its intended functionality. We write the desired
assurance condition as
\[
A_t\models C_t,
\]
where satisfaction denotes the intended property over executions,
rather than a claim that finite testing establishes universal
correctness. In practice, bounded verification and runtime
enforcement provide operational evidence for this condition.

\subsection{What Constitutes a Missing Obligation?}

A new operation without a matching clause is not automatically
a contract gap. It may be irrelevant to $P$, unreachable for
sensitive inputs, or already constrained by another rule.
Conversely, an existing clause may syntactically match an
operation while failing to constrain its policy-relevant effects.
The central question is therefore whether the current contract
permits behavior that violates a supplied safety requirement.

Let $\operatorname{Exec}(\widetilde A_t)$ denote the executions
possible under the candidate agent, and let
$\operatorname{Harm}_p(z)$ indicate that execution $z$ violates
a requirement $p\in P$. An ideal missing-obligation witness
satisfies
\[
\begin{aligned}
z &\in \operatorname{Exec}(\widetilde A_t),\\
z &\models C_{t-1},\\
\operatorname{Harm}_p(z) &\quad\text{for some }p\in P.
\end{aligned}
\]
Here, $z\models C_{t-1}$ means that the execution satisfies
all applicable existing contractual obligations. Thus, the
witness demonstrates a discrepancy between what the current
contract permits and what the fixed policy requires.

To associate the finding with evolution, the execution must
involve a changed operation, a changed guard, or an existing
operation whose reachable inputs or effects are altered by
$\Delta A_t$. This condition focuses analysis on consequences
introduced or exposed by the update, without requiring proof
that the same gap could not have existed previously.

\paragraph{Distinguishing contract gaps from other failures.}
Because ESL uses three-valued predicate evaluation, we distinguish
three outcomes when analyzing a candidate execution.

\begin{itemize}
    \item \textbf{Compliance or enforcement failure.}
    An existing applicable rule evaluates to
    $\mathsf{False}$. The execution violates an existing
    obligation, so adding a new rule is not necessarily required.

    \item \textbf{Observation or grounding failure.}
    An applicable rule evaluates to $\mathsf{Unknown}$
    because relevant objects, bindings, or predicate evidence
    cannot be established. The immediate problem is insufficient
    evidence or instrumentation, rather than a confirmed
    missing obligation.

    \item \textbf{Missing obligation.}
    All applicable existing rules evaluate to
    $\mathsf{True}$, yet the execution violates a supplied
    policy requirement. This provides evidence that the
    current contract does not sufficiently constrain the
    policy-relevant behavior.
\end{itemize}

A confirmed missing-obligation witness therefore requires
definite satisfaction of existing applicable rules, rather
than treating unknown evaluations as successful checks.
When evidence is incomplete, the analysis may seek additional
observations before classifying the behavior.

Our method searches for such witnesses within an explicit
analysis model and resource budget. A discovered witness
demonstrates insufficiency of the current contract for the
observed behavior, but failure to discover one does not
establish completeness. The study evaluates whether the
resulting clauses improve contract compliance and reduce
independently assessed policy violations, while preserving
the agent's ability to perform legitimate tasks. 